\title{Byte Latent Transformer: Patches Scale Better Than Tokens}

\begin{abstract}
We present the Byte Latent Transformer ({\textsc BLT}), a novel byte-level large language model architecture that, for the first time, attains parity with traditional tokenization-based LLMs at scale while delivering notable gains in inference efficiency and robustness. In {\textsc BLT}, individual bytes are aggregated into patches of dynamic size, which constitute the core computational units. These patches are delineated by assessing the entropy associated with upcoming bytes, enabling adaptive allocation of computational resources and model expressiveness to regions of greater data complexity. We offer the first scaling investigation for byte-level models that controls for FLOPs, extending to 8B parameters and 4T training bytes. The findings confirm that models can be successfully scaled when trained directly on byte sequences, obviating the need for a predefined vocabulary. The architecture enhances both training and inference throughput by adaptively forming longer patches in predictable contexts, while also yielding observable improvements in reasoning tasks and out-of-distribution generalization. Ultimately, at equivalent inference costs, {\textsc BLT} achieves markedly superior scaling behavior relative to tokenization-dependent architectures, accomplished through concurrent growth in patch length and model size.
\end{abstract}

\section{Introduction}
\label{section:intro}

We present the Byte Latent Transformer~(\textbf{BLT}), an architecture that dispenses with tokenization by directly learning from raw byte sequences. Remarkably, BLT is the first architecture of its kind to achieve parity with tokenization-based models on large-scale tasks while offering notable gains in terms of efficiency and resilience (\S\ref{sec:robustness}).
Conventional large language models (llms) are predominantly trained in an end-to-end fashion, with the notable exception of tokenization—a heuristic preliminary process that partitions byte sequences into a fixed inventory of tokens. This tokenization process introduces biases during string compression, manifesting as vulnerabilities including sensitivity to domains and modalities~\citep{dagangetting}, susceptibility to input perturbations~(\S\ref{sec:robustness}), insufficient orthographic understanding~\citep{edman2024cute}, and disparities in support for multiple languages~\citep{liang2023xlm,petrov2024language,limisiewicz2024myte}.

Tokenization has historically played a crucial role, as training llms directly on bytes incurs significant computational expenses at scale, mainly because of the extended sequence lengths~\citep{xue2022byt5}. Some studies address this challenge by utilizing more computationally efficient self-attention techniques~\citep{el2020characterbert, clark2022canine} or by leveraging architectures that bypass attention entirely~\citep{wang2024mambabyte}~(\S\ref{section:related_work}). Nonetheless, these strategies chiefly benefit the training of \textit{small models}. For larger models, the bulk of computational overhead in Transformers arises from the extensive feed-forward network layers, which process every byte, while the attention component contributes comparatively little to the overall cost.

We introduce an adaptive and trainable approach for aggregating bytes into \textit{patches}~(\S\ref{section:patching}), accompanied by a novel model architecture that integrates both byte-level and patch-level data. In contrast to tokenization, {\textsc BLT} does not employ a predetermined set of patch vocabulary. Instead, any combination of bytes is projected into latent patch embeddings through efficient, trainable encoder and decoder components. Our results demonstrate that this method yields \textit{greater} computational efficiency compared to models relying on tokenization.

Tokenization-based large language models assign an equal computational budget to each token, irrespective of their individual complexity. This approach prioritizes performance but is often inefficient, as the tokenization heuristics used for data compression do not consistently align with the actual difficulty of generating accurate predictions for each token. Our proposed architecture instead emphasizes adaptive computation, allocating more resources to challenging portions of text as required. For instance, predicting the final characters of most words tends to be straightforward and entails lower entropy, thus typically not necessitating a large transformer, in contrast to the substantially more complex task of predicting the first word of a new sentence. {\textsc BLT}'s design reflects this principle~(\S\ref{section:architecture}) by incorporating a hierarchy of transformer blocks: two lightweight, byte-level \textit{local models}, complemented by a single, more substantial \textit{latent transformer} for capturing broader dependencies~(\autoref{fig:arch_high_level}). 

In order to judiciously allocate computation and determine how bytes should be batched into patches, {\textsc BLT} employs a segmentation strategy guided by the entropy associated with the next-byte prediction. This process results in the creation of byte groupings characterized by relatively consistent information density within their local context.

We introduce the first comprehensive analysis of scaling byte-level models governed by floating point operation (flop) budgets, extending up to 8B parameters and processing as much as 4T training bytes. Our findings demonstrate the feasibility of large-scale, end-to-end training directly from raw bytes, dispensing with the need for fixed-vocabulary tokenization. In particular, {\textsc BLT} achieves training performance on par with Llama 3 when controlled for flops,\footnote{Model computational demands are quantified by tallying the total number of Floating Point OPerations (flops).} but is able to deliver up to 50\% flop reductions during inference~(\S\ref{section:scaling}).

Furthermore, operating directly on raw byte sequences yields notable gains for capturing rare or atypical data patterns. {\textsc BLT} models exhibit superior resilience compared to tokenizer-reliant counterparts when faced with noisy input, and demonstrate stronger proficiency at character-level processing, as shown in their improved performance on orthography, phonological awareness, and low-resource machine translation tasks~(\S\ref{sec:robustness}).

Additionally, with the {\textsc BLT} framework, it is feasible to scale both model and patch sizes concurrently without increasing the inference flop overhead. Employing larger patches typically lowers computational requirements, permitting the surplus to be devoted to expanding the global latent transformer component, since its invocation frequency diminishes. Through inference-flop controlled scaling experiments~(\autoref{fig:fixed_inference_scaling}), we find that our approach achieves far better scaling behaviors relative to models using conventional tokenization-based designs.

In conclusion, this work offers several key contributions: 1) We present {\textsc BLT}, a byte latent llm framework capable of dynamically allocating computational resources to optimize flop efficiency; 2) We demonstrate that our models attain flop-controlled training performance comparable to Llama 3 up to the 8B scale, with the flexibility to achieve up to a 50\% increase in flop efficiency by tolerating slight reductions in evaluation metrics; 3) {\textsc BLT} introduces a novel approach to llm scaling, making it possible to expand model size without exceeding a predetermined inference cost; 4) We provide empirical evidence that {\textsc BLT} models exhibit enhanced robustness against input noise and are sensitive to sub-word structures in the data—properties often overlooked by token-based llms. The full training and inference code for {\textsc BLT} is available at~\url{https://github.com/facebookresearch/blt}.

\section{Patching: From Individual Bytes to Groups of Bytes}
\label{section:patching}

Dividing bytes into discrete \textit{patches} enables {\textsc BLT} to flexibly adjust computational resources according to contextual requirements. As illustrated in Figure~\ref{fig:patching_types}, there are multiple strategies for partitioning bytes into patches. More precisely, a patching function $f_p$ takes a byte sequence $\pmb{x}=\{x_i,|i=1,\ldots n\}$ of length $n$ and divides it into $m < n$ patches $\pmb{p}=\{p_j|j=1,\ldots,m\}$, by assigning each $x_i$ to either 0 or 1, with 1 denoting the beginning of a new patch. For both token-based and patch-based architectures, the computational workload is largely governed by how many steps the primary Transformer must perform; in {\textsc BLT}, this is directly linked to the number of patches generated by the selected patching method. Thus, the typical or average patch size becomes the principal determinant of the computational demands—both at training and inference time—associated with a specific patching approach~(\S\ref{section:flops}).

We now present three patching functions: (1) segmenting into patches of a constant byte count~(\S\ref{section:static-patch}), (2) segmenting at each whitespace character~(\S\ref{section:white-patch}), and (3) dynamic segmentation informed by entropy estimates from a compact byte language model~(\S\ref{section:dyn-patch}). In addition, we address incremental patching and clarify the distinctions between patching and tokenization~(\S\ref{section:bpe-patch}).

\subsection{Strided Patching Every K Bytes}
\label{section:static-patch}

A simple and commonly used method for grouping bytes involves dividing them into fixed-length patches of size $k$, as implemented in MegaByte~\citep{yu2023megabyte}. This approach utilizes a constant stride, which simplifies both the training and inference processes. It also offers an easy way to adjust the typical patch size, thereby allowing for straightforward regulation of computational cost. Nonetheless, this patching strategy introduces notable limitations. Primarily, computational resources are distributed uniformly, rather than in response to the varying informational demands of the data: for instance, valuable transformer computations may be expended on patches containing only whitespace in code, while data-rich regions—such as mathematical notation—may receive insufficient computation. Additionally, this method can result in irregular and context-insensitive segmentation, where identical byte sequences (for example, recurring words) may be partitioned in dissimilar ways.

\subsection{Space Patching}
\label{section:white-patch}

\citet{slagle2024spacebyte} introduces a straightforward enhancement to strided patching by generating new patches after encountering space-like bytes\footnote{
    Space-like bytes refer to any byte not classified as a latin character, digit, or utf-8 continuation byte, and patches are required to include at least one byte that is not space-like.}. Such byte boundaries frequently correspond to linguistic units in many languages. Within the Space patching approach, each word receives a dedicated latent transformer computation (i.e., an increase in flops), ensuring uniformity in how words are segmented across input sequences and focusing computational effort on challenging predictions, which often occur after space characters. For instance, determining the initial byte of a response to “Who composed the Magic Flute?\uline{\hspace{.6em}}” is much more difficult compared to subsequent bytes after “M,” since the first character significantly narrows the plausible answers, rendering the rest of “Mozart” easier to complete. Nevertheless, space patching is limited in its capacity to accommodate all languages and domains, and critically, does not provide adjustable patch sizes. The following section presents a novel patching technique that leverages the observation that the leading bytes in words are generally the most complex to predict, while simultaneously enabling intuitive control over patch size.

\subsection{Entropy Patching: Using Next-Byte Entropies from a Small Byte LM}
\label{section:dyn-patch}

Instead of using rule-based heuristics like whitespace to segment text, we propose a data-driven strategy for detecting regions with high uncertainty in next-byte predictions. Specifically, we present \textit{entropy patching}, a method that defines patch boundaries based on entropy estimations.

We train a small byte-level auto-regressive language model on the training data for {\textsc BLT} and compute next byte entropies under the LM distribution $p_{e}$ over the byte vocabulary $\mathcal{V}$:

\begin{equation}
H(x_i) = \sum_{v \in \mathcal{V}} p_{e}(x_i=v|\pmb{x}_{<i}) \log p_{e}(x_i=v|\pmb{x}_{<i})
\end{equation}

We employ two approaches for detecting patch boundaries based on entropies $H(x_i)$. The first method selects points exceeding a fixed global entropy threshold, as depicted in~\autoref{fig:patching}. The second method highlights points where the entropy is significantly higher compared to the preceding value. Alternatively, this latter strategy can be viewed as pinpointing locations where the approximately monotonic decline of entropy within a patch is disrupted.

\begin{align*}
    \text{Global Constraint}\hspace{1cm}H(x_t)& > \theta_g\\
    \text{Approx. Monotonic Constraint}\hspace{1cm}H(x_t)& - H(x_{t-1}) > \theta_r
\end{align*}

Patch boundaries are determined through a minimal preprocessing process that takes place while loading data. This approach contrasts with~\citet{nawrot-etal-2023-efficient}, who utilize a classifier trained specifically to detect entropy-driven patch boundaries. In our experimental analysis~(\S\ref{section:experiments}), we evaluate these two strategies for discriminating between bytes with low and high entropy.

\subsection{The Byte-Pair Encoding (BPE) Tokenizer and Incremental Patching}
\label{section:incremental-patching}
\label{section:bpe-patch}

A significant number of state-of-the-art llms, such as the baseline Llama 3 considered here, rely on subword tokenization approaches like bpe~\citep{gage1994new, sennrich-etal-2016-neural}. Throughout this work, we use the term ``tokens'' to denote byte-groupings selected from a \textit{finite} vocabulary that is established before training begins, in contrast to ``patches'', which denote adaptively formed groups of sequences that are not constrained to a predetermined vocabulary. One key distinction between these approaches is that with tokens, the model is unable to directly leverage the raw byte features.

A key advancement offered by {\textsc BLT} in comparison to tokenization-based models lies in its redefinition of the relationship between vocabulary size and computational requirements. In conventional large language models, expanding the vocabulary leads to an increase in the average token length, resulting in fewer decoding steps; however, this also enlarges the output dimension of the model’s final projection layer. This inherent compromise restricts the degree to which tokenization-based strategies can significantly alter both token sizes and the associated inference overhead. As an illustration, Llama 3 achieves a boost in average token length—from 3.7 to 4.4 bytes—but only by expanding the capacity of its embedding matrix by a factor of four relative to Llama 2.

During sequence generation, {\textsc BLT} must determine if the current position in the byte sequence coincides with a patch boundary, as this influences whether additional computation through the Latent Transformer is necessary. This assessment must be made without relying on any bytes that have not yet been produced, meaning the patching strategy cannot depend on knowledge of future parts of the sequence to segment the bytes. More formally, a patching function $f_p$ is called incremental if it fulfills the following condition:
$$f_p(\pmb{x}_{<i}) = f_p(\pmb{x})_{<i}$$
The bpe patching method does not qualify as incremental because the encoding of a given prefix may vary depending on how the sequence continues, and thus fails to meet the criterion above\footnote{Appending a dedicated delimiter token to mark patch boundaries can render bpe incremental, albeit at the cost of lengthening the byte sequence.}.

\section{{\textsc BLT} Architecture}
\label{section:architecture}
The {\textsc BLT} framework consists of a substantial global autoregressive language model responsible for processing patch representations. Accompanying this global component are two compact local models: one tasked with transforming sequences of bytes into patch representations, and another dedicated to reconstructing bytes from these patch representations~(\autoref{fig:arch_high_level}).

\subsection{Latent Global Transformer Model}

\textit{The Latent Global Transformer} refers to an autoregressive transformer architecture, denoted as $\mathcal{G}$, that comprises $l_{\mathcal{G}}$ layers. This model processes an input sequence of latent patch embeddings, $p_j$, and produces a corresponding output sequence of patch embeddings, $o_j$. In our notation, the index $j$ consistently designates individual patches, whereas $i$ is reserved for indexing bytes. The attention mechanism employed by the global model is block-causal~\citep{dubey2024llama}, ensuring that each patch's attention is constrained to its own and preceding patches within the given document. Notably, the global model is responsible for the majority of the floating point operations required during both pre-training and inference. Accordingly, by selectively activating this model, we gain the capacity to adaptively manage computational costs across different regions of the input and output, reflecting variations in their respective complexities.

\subsection{Local Encoder}
\label{section:local}

\textit{The Local Encoder Model}, represented by $\mathcal{E}$, constitutes a compact transformer-based architecture comprising $l_{\mathcal{E}}$ layers, with $l_{\mathcal{E}}$ significantly smaller than $l_{\mathcal{G}}$. Its principal function is to rapidly project input byte sequences $b_i$ into rich patch-level embeddings, denoted as $p_j$. Notably diverging from standard transformers, the model incorporates a cross-attention layer subsequent to every transformer layer; this modification serves to aggregate information from the byte-level representations to form patch representations~(\autoref{fig:crossattn}).

At the initial stage, individual bytes $b_i$ are mapped into the embedding space using a $\mathbb{R}^{256 \times h_{\mathcal{E}}}$ embedding matrix, yielding vectors $x_i$. There is also an option to enhance these embeddings with hash-based embeddings as explained in~(\S\ref{sec:hash-emb}).

Through an alternating stack of transformer and cross-attention modules~(\S\ref{sec:enc-cross-attn}), the byte representations are progressively converted into patch embeddings, $p_i$, which are then forwarded for further processing by the global transformer $\mathcal{G}$.

The transformer modules adopt a \textit{local block causal} attention scheme. Within this mechanism, each byte is permitted to attend to the previous $w_{\mathcal{E}}$ bytes within a fixed window. While this window may extend across patch boundaries (which are determined dynamically), it does not traverse the boundaries of distinct documents. Details concerning both the embedding techniques and the internal structure of the cross-attention block are provided in the following sections.

\subsubsection{Encoder Hash n-gram Embeddings}
\label{sec:ngram-emb}
\label{sec:hash-emb}
An essential aspect of building resilient and rich representations at every position $i$ is the integration of contextual data from the bytes that came before. In {\textsc BLT}, this is accomplished by representing the current byte $b_i$ both as an isolated element \textit{and} within a sequence of bytes—namely, a byte n-gram. At every step $i$, we begin by assembling the corresponding byte-grams

\begin{align}
    g_{i,n}=\{b_{i-n + 1},\ldots, b_i\}
\end{align}

for each byte index $i$ and for $n$ values ranging from three to eight.\footnote{
We do not consider byte-grams of length $n$ or greater when $i<n$. }

Subsequently, we propose the use of hash $n$-gram embeddings. This method applies a hash function to each byte $n$-gram to determine its corresponding index in the embedding matrix $E_{n}^{hash}$, which remains of constant size for every $n \in\{3,4,5,6,7,8\}$~\citep{bai2010learning}. The computed embedding is then summed with the embedding of the current byte, after which it undergoes normalization before being supplied as the input to the local encoder architecture. The resulting combined embedding is expressed as:

\begin{align}
e_i &= x_i + \sum_{n = 3,...,8}E_{n}^{hash}(\text{Hash}(g_{i,n})) \\
\text{where, Hash}(g_{i,n}) &= \text{RollPolyHash}(g_{i,n}) \% |E_{n}^{hash}|
\end{align}

Each $e_i$ is scaled by dividing by the sum of the different $n$-gram sizes incremented by one, employing RollPolyHash as outlined in Appendix~\ref{appendix:rollpolyhash}. Section~\ref{section:ablations} presents an analysis that isolates the impact of $n$-gram hash embeddings, specifically examining variations in $n$ and the size of the embedding table with regard to flop-controlled scaling law behavior. Besides investigating hash-based $n$-gram embeddings, we also evaluated frequency-based $n$-gram embeddings; further information on this line of inquiry is available in Appendix~\ref{appendix:ngrams}.

\subsubsection{Encoder Multi-Headed Cross-Attention}
\label{sec:enc-cross-attn}
Our approach largely adopts the cross-attention input mechanism from the Perceiver architecture~\citep{jaegle2021perceiver}; however, the principal distinction lies in our use of latent representations that dynamically correspond to individual patch representations rather than a static collection of latent vectors~(\autoref{fig:crossattn}). Additionally, each latent is restricted to attending solely to the bytes within its own associated patch. In this setup, each patch $p_j$ is linked to a query vector, which is obtained by aggregating the byte-level representations for patch $p_j$ through a pooling operation, followed by projection via a linear transformation $\mathcal{E}_{C} \in \mathbb{R}^{h_{\mathcal{E}} \times (h_{\mathcal{E}} \times U_{\mathcal{E}})}$, where $U_{\mathcal{E}}$ designates the encoder’s number of cross-attention heads. More precisely, letting $f_{\text{bytes}}(p_j)$ represent the sequence of bytes for a given patch $p_j$, the calculation proceeds as follows

\begin{align}
P_{0,j} &= \mathcal{E}_{C}(f_\text{bytes}((p_j)), f~ \text{is a pooling function} \\
P_l &= P_{l-1} + W_o\left(\text{softmax}\left(\frac{QK^T}{\sqrt{d_k}}\right)V\right) \\
\text{where } Q_j &= W_q(P_{l-1,j}), K_i = W_k(h_{l-1,i}), V_i = W_v(h_{l-1,i}) \\
h_l &= \text{Encoder-Transformer-Layer}_l(h_{l-1})
\end{align}

Here, $P \in \mathbb{R}^{n_p \times h_{\mathcal{G}}}$ denotes the $n_p$ representations of patches that are input to the global model, initialized via a pooling operation over the byte embeddings $e_i$ linked to each patch $p_j$. The parameters $W_q$, $W_k$, $W_v$, and $W_o$ serve as projection matrices for queries, keys, values, and outputs, respectively, where keys and values result from applying these projections to the byte-level representations $h_i$ from the preceding layer (or $e_i$ if in the first layer). For attention computation, a patch-specific masking approach is adopted: each query $Q_j$ only attends to the keys and values originating from bytes in its own patch $j$. With multi-headed attention applied to $Q$, $K$, and $V$, and noting that patch-level representations often have higher dimensionality ($h_{\mathcal{G}}$) than $h_{\mathcal{E}}$, we represent $P_l$ using several heads each of size $h_{\mathcal{E}}$ during cross-attention, concatenating these heads to yield representations in $h_{\mathcal{G}}$ dimensions. The module employs a pre-LayerNorm on the queries, keys, and values and omits the use of positional embeddings in the cross-attention operation. The cross-attention block is finally wrapped with a residual connection.

\subsection{Local Decoder} 

The local decoder $\mathcal{D}$ mirrors the structure of the local encoder, functioning as a compact transformer-based architecture that uses $l_{\mathcal{D}}$ layers, where $l_{\mathcal{D}} << l_{\mathcal{G}}$. This decoder maps a set of global patch representations $o_j$ into a sequence of raw bytes, $y_i$. By conditioning each predicted byte on previously decoded bytes, the local decoder incorporates the local encoder’s hidden states corresponding to the byte sequence as inputs. It consists of $l_{\mathcal{D}}$ repeated blocks, each containing a cross-attention layer followed by a transformer layer. The decoder's cross-attention module operates first, transforming patch-level representations into byte-level representations, after which the standard transformer block further processes these byte-wise representations.

\subsubsection{Decoder Multi-headed Cross-Attention} 

In the decoder cross-attention, the roles of the queries and key/values are interchanged i.e. the byte-representations are now the queries, and the patch representations are now the key/values. The initial byte-representations for the cross-attention are initialized as the byte embeddings from the last encoder layer i.e. $h_{l_{\mathcal{E}}}$. The subsequent byte-representations for layer $l$, $d_{l,i}$ are computed as:

\begin{align}
D_0 &= h_{l_{\mathcal{E}}} \\
B_l &= D_{l-1} + W_o\left(\text{softmax}\left(\frac{QK^T}{\sqrt{d_k}}\right)V\right), \\
  &\text{where } Q_i = W_q(d_{l-1,i}), K_i = W_k(\mathcal{D}_C(o_j)), V_i = W_v(\mathcal{D}_C(o_j)) \\
  D_l &= \text{Decoder-Transformer-layer}_l(B_l)
\end{align}

In this context, $W_k$ and $W_v$ denote the key and value projection matrices responsible for executing a linear transformation and a subsequent split operation $\mathcal{D}_C$ on the terminal patch representations $o_j$ generated by the global model. $W_q$ serves as the query projection matrix acting on the byte representations $d_{l-1}$ stemming from the prior layer in the decoder transformer (or $h_{l_{\mathcal{E}}}$ in the case of the initial layer). $W_o$ functions as the output projection matrix. Consequently, this formulation yields $B \in \mathbb{R}^{h_{\mathcal{D}} \times n_b}$, with $n_b$ indicating the count of output bytes. The subsequent set of decoder representations, $D_l$, are derived by applying a decoder transformer layer to the outputs from the cross-attention module, $B$. This procedure mirrors the approach in the local encoder cross-attention by adopting multi-head attention, employing pre LayerNorms, omitting positional embeddings, and including a residual connection around the cross-attention structure.

\section{Experimental Setup}
\label{section:experiments}
We construct a series of methodically controlled experiments to assess {\textsc BLT} in comparison to models based on tokenization, taking deliberate care to prevent {\textsc BLT} from benefitting from any potential sequence length advantages.

\subsection{Pre-training Datasets} In this work, all model variants under consideration are initially trained using two primary datasets: 1) the Llama 2 dataset~\citep{touvron2023llama}, encompassing 2 trillion tokens sourced from diverse publicly available materials that undergo rigorous cleaning and filtering processes to ensure higher quality; and 2) {\textsc BLT}-1T, a newly assembled dataset containing 1 trillion tokens, drawn from assorted public sources and including a portion of the pre-training corpus made available by Datacomp-LM~\citep{li2024datacomp}. The Llama 2 dataset is employed for scaling law analyses, specifically utilizing the recommended token counts suggested by~\cite{dubey2024llama}, in order to optimize the architectural design of {\textsc BLT}. On the other hand, the {\textsc BLT}-1T dataset is used to perform a full pre-training cycle, facilitating direct downstream comparison with Llama 3. It is important to emphasize that neither dataset incorporates any content from Meta’s products or services. For our baseline models using tokenizers, we employ the Llama 3 tokenizer with a 128K-token vocabulary, as this configuration yielded better baseline results compared to the Llama 2 tokenizer during our evaluations.

\subsection{Entropy Model}  
For our experiments, the entropy model utilized is a byte-level language model, trained using the identical data distribution as the comprehensive {\textsc BLT} model. By default, this model is a transformer architecture with 100 million parameters, organized across 14 layers and featuring a hidden size of 512. It employs sliding window attention over 512 bytes. Other hyperparameters mirror those chosen for both our local and global transformer models. We also explored variations in the model's size, receptive field, and architectural design, detailed in \autoref{section:ablations}. Notably, when configured with a sufficiently limited receptive field, the trained entropy model may be efficiently represented as a lookup table.

\subsection{Entropy Threshold and Equalizing Context Length} For entropy-based patching approaches, we determine a patching threshold aimed at yielding a specific average \textit{patch size} across the pretraining data mixture. In {\textsc BLT}, the \textit{patch size}, unlike that derived through standard tokenization, is a flexible parameter that can be set arbitrarily, thereby strongly affecting the amount of context the model processes. To ensure that models do not benefit unfairly from increased patch size—through exposure to more context—we regulate the setup such that the expected number of bytes per batch stays constant. Accordingly, sequence lengths are shortened for models that utilize larger patch sizes. For models trained on Llama 2 data, we fix the context window to 8k bytes, whereas for those trained on the {\textsc BLT}-1T dataset, the context window is extended to 16k bytes on average, with the batch size held steady at an average of 16M bytes.

Although the mean batch size remains unchanged, dynamic patching approaches result in varying proportions of bytes per patch when assembling data batches. To enhance efficiency, our {\textsc BLT} training approach composes batches based on patch counts, thereby eliminating the need for padding within the more computationally intensive latent transformer module. This method guarantees that each batch includes an identical count of patches. In the training process, byte sequences are either padded or truncated—set to 12k bytes for Llama 2 and 24k bytes for {\textsc BLT}-1T datasets—to prevent unexpected increases in memory usage caused by exceptionally large patches.

\subsection{Entropy Model Context}
\label{section:entropy-context}
Our experiments reveal that employing entropy patching tends to result in larger patches, particularly within highly structured domains such as multiple choice tasks, where content often exhibits a high degree of repetition (refer to patching in an MMLU example in \autoref{fig:mmlu_patching}). These increases in patch size stem from the entropy model context assigning lower entropy values to repeated sequences. For the comprehensive execution of {\textsc BLT}-Entropy utilizing a patch size of 4.5, we mitigate this by resetting the entropy context with the insertion of new lines and by adopting an approximate monotonicity constraint, which is less prone to "entropy drift" due to fluctuations in context length. Notably, this adjustment solely pertains to the method of entropy calculation, leaving the overall approach for determining the entropy threshold unchanged.

\subsection{FLOPs Estimation} 
\label{section:flops}

Our methodology closely adheres to the approach described in Chinchilla~\citep{hoffmann2022training} for calculating the floating-point operations (flops) in transformers, which encompasses contributions from the feed-forward components, the qkvo projections within self-attention, as well as the attention score calculation and output projection. However, our analysis diverges by modeling the input embedding layer as a highly efficient lookup table rather than a conventional dense matrix multiplication, effectively reducing its flop count to zero. Consistent with established practice, we approximate the computational load of the backwards pass as being twice that of the forward pass.

To compute flops \textit{per byte} for {\textsc BLT} models, we add up the flops for the local encoder transformer, the global latent transformer, and the local decoder transformer, together with the cross attention blocks in the encoder and the decoder:

\begin{align}
    \text{FL}_{\text{{\textsc BLT}}} &= \frac{1}{n_p} \text{Transf. FL}(h_{\mathcal{G}}, l_{\mathcal{G}}, m=n_{ctx}/n_p, V=0)\\
   &\quad+ \text{Transf. FL}(h_{\mathcal{E}}, l_{\mathcal{E}}, m=w_{\mathcal{E}}, V=0) \\
   &\quad+ \text{Transf. FL}(h_{\mathcal{D}}, l_{\mathcal{D}}, m=w_{\mathcal{D}}, V=256) \\ 
    &\quad+ \text{Cross Attn. FL}(h_{\mathcal{E}}, l_{\mathcal{E}}, m=n_p, r=n_p/k) \cdot \frac{k}{n_p} \\ 
   &\quad+ \text{Cross Attn. FL}(h_{\mathcal{D}}, l_{\mathcal{D}}, m=k, r=k/n_p)
\end{align}

Here, $n_{ctx}$ denotes the length of the input sequence measured in bytes, $n_p$ refers to the patch size, and $r$ indicates the proportion of queries to key/values. The parameter $k$ defines the ratio between patch-dimension and byte-dimension, representing the count of local model partitions concatenated to compose a global model embedding ($k=2$ as illustrated in \autoref{fig:crossattn}). $V$ signifies the size of the output projection vocabulary, relevant solely for the local decoder. The context length $m$ in the attention mechanism differs depending on whether the operation is carried out over the byte sequence or the patch sequence. We adjust the corresponding computation of attention-related flops for each module as appropriate. Appendix~\ref{section:flops-appendix} details the precise formulae for calculating flops in both Transformer and Cross-Attention components.

\subsection{Bits-Per-Byte Estimation} Perplexity only makes sense in the context of a fixed tokenizer as it is a measure of the uncertainty for each token. When comparing byte and token-level models, following previous work~\citep{xue2022byt5, yu2023megabyte, wang2024mambabyte}, we instead report Bits-Per-Byte (BPB), a tokenizer independent version of perplexity. Specifically:

\begin{align}
    \text{BPB}(x)&= \frac{\mathcal{L}_{CE}(\pmb{x})}{\ln(2)\cdot n_{\text{bytes}}}
\end{align}

Here, the aggregate cross-entropy loss, reflecting the uncertainty associated with the data $\pmb{x}$, is divided by both the total number of bytes in $\pmb{x}$ and a constant factor for normalization.

\subsection{Transformer Architecture Hyperparameters} 

Throughout all transformer blocks within {\textsc BLT}—encompassing both the local and global modules—we adopt most of the architectural components established in Llama~3~\citep{dubey2024llama}. The feed-forward networks utilize the SwiGLU activation function~\citep{swiglu}, and self-attention mechanisms incorporate rotary positional embeddings (RoPE)~\citep{RoPe} parameterized with $\theta=500000$~\citep{xiong2024effective}. Layer normalization is achieved using RMSNorm~\citep{RMSNorm}. For self-attention layers employing standard attention masks, such as \textit{block causal} and \textit{fixed-window block causal}, we rely on Flash Attention~\citep{dao2022flashattention}; for fixed-width masks, the window size is set to 512. In cross-attention layers, where attention masks depend dynamically on patch selection, we employ Flex Attention\footnote{\url{https://pytorch.org/blog/flexattention}}, which enables fused operations and considerably accelerates the training process.

\subsection{{\textsc BLT}-Specific Hyperparameters} In order to evaluate the performance of the {\textsc BLT} models, we design experiments targeting two main aspects: scaling behavior and performance on downstream tasks. We analyze models with varying sizes—400M, 1B, 2B, 4B, and 8B parameters—for these studies. Detailed architectural hyperparameter configurations can be found in Appendix~Table~\ref{tab:arch}.  
For query initialization in the initial cross-attention layer of the local encoder, we employ max-pooling. Across all {\textsc BLT} variants, we implement $500,000$ hashes using a single hash function, considering n-gram windows between 3 and 8. All models share a fixed learning rate of $4e-4$, as hyperparameter exploration over the $1e-3$ to $1e-4$ range at 400M and 1B model sizes established that identical learning rates are optimal for both token and {\textsc BLT} variants. When assessing scaling trends on the Llama-2 dataset, we adopt the batch-size recommendations (or their byte-equivalent) put forth by~\cite{dubey2024llama}. For optimization, we leverage the AdamW optimizer~\citep{adamw}, assigning $\beta_1=0.9$ and $\beta_2=0.95$, with $\epsilon=10^{-8}$. The learning rate follows a linear warm-up over 2000 steps, then decays to zero via a cosine annealing schedule. Additionally, a weight decay rate of 0.1 is imposed and gradient norms are globally clipped at 1.0.

\section{Scaling Trends}
\label{section:scaling}

We offer a comprehensive analysis of the scaling behaviors exhibited by byte-level models, providing guidance for the future scaling of {\textsc BLT} models. Our scaling investigation addresses certain gaps in existing research on byte-level approaches in the following respects: (a) We examine scaling patterns under compute-optimal training conditions; (b) We construct and assess 8B-parameter models trained with substantial datasets (up to 1T tokens or 4T bytes), evaluating their effectiveness on a range of downstream benchmarks; and (c) We assess scaling effects when inference cost is held constant. Subsequent sections will delve into the unique benefits derived from modeling sequences at the byte level.

\subsection{Parameter Matched Compute Optimal Scaling Trends}
\label{section:parameter-matched-scaling}
We conduct experiments on the Llama 2 dataset by training multiple \textit{compute-optimal} bpe and {\textsc BLT} models at four distinct model scales, spanning 1B to 8B parameters. Training flops are then graphed against language modeling accuracy using a representative sample from the overall training data mixture. For the bpe models, we apply the parameter-to-data ratio identified as optimal in Llama~3~\citep{dubey2024llama}, aligning our procedure with established \textit{compute-optimal} practice to yield maximal training data performance within a fixed compute allocation~\citep{hoffmann2022training}. This serves as a strong reference point for comparing model types. Corresponding to each bpe model, a {\textsc BLT} model is trained on the identical dataset, utilizing a Latent Transformer whose scale and architecture mirrors that of the matched bpe Transformer counterpart.

As depicted in~\autoref{fig:scaling-compute-opt} (right), {\textsc BLT} models consistently equal or surpass the performance of their bpe equivalents, a pattern that persists as both model size and computational resources increase. To our knowledge, {\textsc BLT} represents the first byte-level Transformer to reach parity in scaling performance with BPE-based architectures under compute-optimal conditions. This supports our hypothesis that the ideal proportion of parameters to training compute established for bpe is similarly applicable to {\textsc BLT}, or differs only minimally.

Both enhancements in model architecture and the implementation of dynamic patching play essential roles in aligning with bpe scaling patterns. In ~\autoref{fig:scaling-compute-opt} (left), we present a comparison between models using space-patching techniques and Llama 3. SpaceByte \citep{slagle2024spacebyte} is approximated by employing {\textsc BLT} with space-patching, but omitting n-gram embeddings and cross-attention mechanisms. While SpaceByte demonstrates gains relative to Megabyte, it does not yet reach the performance of Llama 3. As shown in \autoref{fig:scaling-compute-opt} (right), we visualize how both the architectural advancements and dynamic patching contribute to model performance. Notably, at large scales, {\textsc BLT} models match the performance of leading tokenizer-based systems such as Llama 3.

Additionally, our findings indicate that for tokenizer-based models, selecting the Llama-3 tokenizer leads to improved performance compared to using the Llama-2 tokenizer, even when both are trained on identical data.

Ultimately, our {\textsc BLT} framework exhibits performance characteristics that lie between Llama 2 and Llama 3 when employing substantially larger patch sizes. Whereas the bpe tokenizers in Llama 2 and 3 encode tokens with average lengths of 3.7 and 4.4 bytes respectively, {\textsc BLT} attains analogous scaling behavior even when the mean patch size increases to 6 or 8 bytes. Since inference flops decrease as average patch size grows, choosing an 8-byte patch can reduce inference computational demands by approximately 50\%. Furthermore, as the size of both the model and dataset expands, models trained with larger patches tend to exhibit superior performance. While {\textsc BLT} initialized with an 8-byte patch size initially lags behind bpe Llama 2 at the 1B scale, it outperforms its bpe counterpart by the time the model reaches 7B scale. This indicates the potential of such large patch sizes to yield greater benefits as models and computational budgets continue to increase, and suggests that even larger patch sizes may become practical for future large-scale training regimes.

\subsection{Beyond Compute Optimal Task Evaluations}
\label{section:evals}
To further examine how scaling affects performance, we extend training for an 8B {\textsc BLT} model beyond the previously identified compute optimal regime, utilizing the expanded and higher-quality {\textsc BLT}-1T dataset. We evaluate this model on a set of established classification and generation benchmarks. Our evaluation focuses on tasks involving common sense reasoning, general world knowledge, and code generation:
\paragraph{Classification tasks} cover ARC-Easy (0-shot)~\citep{Eval_ARC}, Arc-Challenge (0-shot)~\citep{Eval_ARC}, HellaSwag (0-shot)~\citep{Eval_hellaswag}, PIQA (0-shot)~\citep{Eval_piqa}, and MMLU (5-shot)~\citep{hendrycks2020measuring}. For these tasks, we adopt a prompt-scoring approach, where choice character likelihoods are used, and the mean accuracy is reported.
\paragraph{Coding related generation tasks:}  The capability of LLMs to produce Python code is assessed through pass@1 results for MBPP (3-shot)~\citep{Eval_MBPP} and HumanEval (0-shot)~\citep{Eval_HumanEval}.

In \autoref{tab:evals}, we present a comparison among three models trained on the {\textsc BLT}-1T dataset: a Llama 3 model utilizing a bpe tokenizer,\footnote{We opt for the Llama 3 tokenizer with its 128k vocabulary, as it demonstrates superior performance compared to Llama 2's 32k vocabulary.} alongside two distinct versions of the {\textsc BLT} architecture. One variant implements a space-patching mechanism ({\textsc BLT}-Space), while the other applies an entropy-based patching approach ({\textsc BLT}-Entropy), both subject to an approximate monotonicity constraint. Additionally, for the entropy-based model, we reset the model's context using new lines as detailed in~\autoref{section:entropy-context}. Each of the three models is trained using a similar computational (flop) budget. Notably, in the case of {\textsc BLT}-Entropy, we further adjust the entropy threshold at inference from 0.6 to 0.1, an alteration that yields enhanced task outcomes but results in an increased number of inference steps.

The {\textsc BLT}-Entropy model surpasses the Llama 3 model in 4 of the 7 evaluated tasks, despite both models being trained on an equal byte count. This superior performance is likely attributable to two main factors: (1) more effective utilization of computational resources enabled by dynamic patching, and (2) the explicit modeling of information at the byte level rather than at the token level.

Conversely, {\textsc BLT}-Space demonstrates inferior performance to the Llama 3 tokenizer across all tasks except one. However, it offers a notable decrease in inference flops, attributed to its larger average patch size of 6 bytes. In contrast, models utilizing bpe and entropy-patching strategies maintain an average patch size near 4.5 bytes when applied to the training dataset mixture. Given an identical training budget, the model with the larger patch size processes 30\% more data than the alternative models, a factor that may move {\textsc BLT} further from achieving compute optimality.

\subsection{Patches Scale Better Than Tokens} {\textsc BLT} architectures allow for concurrent scaling of both the model and patch sizes, without exceeding a fixed training and inference compute budget or altering the volume of training data. Unlike conventional token-based models bound by a predetermined vocabulary, patch-based models uniquely permit arbitrary enlargement of patch size, circumventing the traditional efficiency constraints, as outlined in Section~\ref{section:bpe-patch}. Utilizing larger patches reduces computational requirements, which enables additional capacity to be devoted to enlarging the global latent transformer, since it needs to process information less frequently.

A fixed inference scaling experiment is performed to examine whether scaling up model size—while using larger patch sizes and fewer inference steps—outperforms using smaller models that require more steps. Utilizing starting points of 400 million and 3.6 billion parameters with the Llama 2 tokenizer, we select flop-matched counterparts employing the Llama 3 tokenizer, as well as {\textsc BLT}-Entropy variants that use average patch sizes of 6 and 8 bytes on the mixed training data (see~\autoref{tab:fixed-inf-params} for model specifications). In the case of models with a patch size of 8, the architecture includes 3 encoder layers, in contrast to just 1 for others. All models are trained under several different flop-constrained training regimes.

\autoref{fig:fixed_inference_scaling} illustrates that {\textsc BLT} models exhibit more favorable scaling behavior compared to tokenization-based models across both classes of inference flop constraints. For both settings, BPE-based models initially outperform at lower training compute levels, yet {\textsc BLT} overtakes them shortly after the point of compute-optimality. In real-world scenarios, it is often advantageous to allocate greater resources during the initial pretraining phase, resulting in models that deliver superior performance within a set inference compute limit. The 8B parameter model class, such as Llama 3.1, exemplifies this approach, having been trained on data volumes that exceed the compute-optimal amount for that model scale by two orders of magnitude.

The point at which {\textsc BLT} surpasses token-based architectures has moved marginally closer to the compute-optimal regime with the transition to models in the higher flop class, decreasing from approximately 3x to 2.5x the compute-optimal allocation. Moreover, models employing a patch size of 8 demonstrate a sharper scaling rate in this higher flop regime, allowing them to outperform other models at an earlier stage. As elaborated in Section~\ref{section:parameter-matched-scaling}, models with greater patch sizes tend to approach the performance of BPE models when scaled up. This phenomenon can be partly explained by the fact that, as model size increases, the proportion of total flops attributed to the byte-level Encoder and Decoder decreases, given their slower scaling relative to the Latent Transformer. When the model’s total parameters expand by a factor of 20, from 400M to 8B, only a modest increase—roughly a doubling—occurs in {\textsc BLT}’s local parameters. Notably, increases in patch size directly impact flops in the patch Latent Transformer rather than the byte-level components. Consequently, this accounts for the observed increase in the size ratio of {\textsc BLT}-Entropy ps=8 relative to Llama 2, growing from 1.6x to 1.7x as the model size was scaled upward.

To summarize, our analysis on patch-length scaling reveals that the {\textsc BLT} architecture benefits from concurrent increases in both patch length and model capacity, leading to superior scaling behavior. Notably, these advantageous scaling patterns appear to continue, and even strengthen, as model size is further increased.

\section{Byte Modeling Improves Robustness} In addition, we evaluate the resilience of {\textsc BLT} in contrast with models based on tokens that do not inherently leverage byte-level data, and we introduce a strategy for converting pretrained token-based models to operate at the byte level.
\label{sec:robustness}

\subsection{Character-Level Tasks}

One of the initial reasons for developing byte-level models was their ability to handle input noise at the byte level, as well as their inherent understanding of the sub-components of tokens—a capability that tokenizer-based models often lack. To assess these characteristics, we conduct further experiments using benchmarks that test both the resilience to noisy input and the recognition of individual characters, encompassing English and multilingual contexts, along with digits and phonemes. The outcomes of these evaluations are summarized in Table~\ref{tab:char_tasks}.

\paragraph{Noisy Data} To assess and compare the robustness of tokenizer-based models against that of {\textsc BLT}, we generate corrupted versions of the benchmark classification datasets referenced in Section~\ref{section:evals}. Specifically, we utilize five unique character-level perturbation techniques to systematically modify the textual data: (a) \textit{AntSpeak}: transforms the text so that every character is capitalized and separated by spaces; (b) \textit{Drop}: randomly deletes 10\% of the characters from the text; (c) \textit{RandomCase}: assigns random casing such that half the characters are in uppercase and the other half in lowercase; (d) \textit{Repeat}: selects 20\% of the characters and repeats them as many as four consecutive times; (e) \textit{UpperCase}: converts all alphabetic characters to uppercase. For evaluation purposes, each of these noising strategies is independently applied to either the prompt, the completion, or both, constituting distinct experimental settings, and mean performance metrics are subsequently reported. Results for the perturbed HellaSwag~\citep{Eval_hellaswag} dataset are summarized in Table \ref{tab:char_tasks}, where {\textsc BLT} demonstrates consistently greater resilience to noise compared to tokenizer-based approaches—achieving an average performance improvement of 8 points over models trained on identical data, and even outperforming the Llama 3.1 model, which was trained on a considerably larger dataset.

\paragraph{Phonology - Grapheme-to-Phoneme (G2P)} We evaluate how effectively {\textsc BLT} can convert sequences of graphemes, which are the characters that make up a word, into their corresponding phonemic transcriptions. Table \ref{tab:char_tasks} summarizes the performance on the G2P task under a 5-shot learning scenario with Phonology Bench~\citep{suvarna2024phonologybench}. Our analysis demonstrates that {\textsc BLT} surpasses the baseline model based on the Llama 3 1T tokenizer in this task.

\paragraph{CUTE}
To evaluate character-level comprehension, we test {\textsc BLT} on the CUTE benchmark~\citep{edman2024cute}, which features a suite of tasks grouped into three primary types: evaluating compositionality, assessing orthographic resemblance, and measuring the capacity for sequence manipulation. This benchmark is notably challenging for conventional tokenizer-based models, which tend to recognize the spelling of their tokens but encounter difficulties leveraging this information for text manipulation tasks. As presented in Table~\ref{tab:char_tasks}, {\textsc BLT}-Entropy exceeds the performance of both BPE-based Llama 3 models by over 25 points on the CUTE suite. Our model, in particular, excels in the manipulation of characters, attaining an accuracy of 99.9\% on both spelling-focused tasks. Remarkably, these substantial gains are achieved despite {\textsc BLT} being trained on 16 times less data compared to Llama 3.1, suggesting that capturing character-level representations poses a significant challenge for BPE architectures. Figure \ref{fig:cute_examples} provides examples where Llama 3 with a tokenizer underperforms, in contrast to our {\textsc BLT} model’s robust outcomes. BPE outperforms only on tasks involving word deletion and insertion; such manipulations are less direct for a byte-level model, yet the margin is not substantial, and building up from characters to words may prove less complex than the reverse process. All experiments employ consistent evaluation protocols and utilize original prompts from Huggingface. It is also possible that BPE-based models could show improved results with more extensive prompt engineering.

\paragraph{Low Resource Machine Translation} We assess the performance of {\textsc BLT} for both translation to and from English across six major language families, focusing on twenty-one lower resource languages representing a range of scripts, as included in the FLORES-101 benchmark~\citep{goyal2022flores}. The results, summarized in Table~\ref{tab:flores} using SentencePiece BLEU scores, indicate that {\textsc BLT} consistently surpasses a model utilizing the Llama 3 tokenizer, achieving a 2-point overall gain when translating into English and a 0.5-point improvement in translations from English. While {\textsc BLT} matches or marginally exceeds Llama 3's performance on popular language pairs, it displays even greater strength for many language pairs from less-resourced families. These findings highlight byte modeling’s advantage in extending model generalization to rare and complex byte sequences.

\subsection{Training {\textsc BLT} from Llama 3}
\label{sec:distillation}
We investigate a methodology in which {\textsc BLT} models benefit from established pre-trained, tokenizer-based architectures to improve and accelerate their learning dynamics. This is accomplished by initializing the global transformer parameters within {\textsc BLT} with those from a Llama 3.1 model that has already been pre-trained. In the subsequent training phase, the global transformer’s weights are updated using a learning rate set to one-tenth of that used for updating both the local encoder and local decoder components. The number of training steps matches the recommended setting for Llama 3. The results, as shown in Table~\ref{tab:distill}, indicate that initializing {\textsc BLT} from Llama 3.1 yields clear advantages over both Llama 3 and a baseline {\textsc BLT}—all under identical flop budgets. Furthermore, a comparison with our {\textsc BLT}-Entropy variant (refer to Table~\ref{tab:evals}), which was exposed to a much larger corpus (1T tokens), reveals that the {\textsc BLT} initialized from Llama 3.1 still surpasses it on the MMLU benchmark. This demonstrates the potential of this transfer approach in achieving significant reductions in required training compute.

This approach can alternatively be interpreted as transforming models reliant on tokenizers into ones that operate without them, thereby adapting a pre-trained LLaMA 3.1 model into the {\textsc BLT} framework. For a thorough assessment, we report results for both the original LLaMA 3.1—trained on 15T tokens—and its corresponding {\textsc BLT} variant in Table \ref{tab:distill}. While the adapted model shows only minor decreases in performance on MMLU and HumanEval, there is a more pronounced decline across other benchmarks. These findings indicate that additional research is necessary to make full use of the pre-trained model’s capabilities, particularly by refining data compositions and tuning hyperparameters.

\section{Ablations and Discussion}
\label{section:ablations}
Here, we present ablation studies that support our selection of architecture for {\textsc BLT}, as well as the decisions regarding the patching approach and hyper-parameter values for the {\textsc BLT} model with 8B parameters trained on the {\textsc BLT}-1T dataset.

\paragraph{Entropy Model Hyper-parameters} To investigate how the size of the entropy model and the length of the context window influence scaling behavior, we experiment with byte-level entropy transformer architectures, systematically varying the parameter count from 1m to 100m and the context window length between 64 and 512 tokens. For analysis, we generate bpb versus training FLOP scaling curves by leveraging our $400m$ and $1b$ {\textsc BLT} models, each trained on the Llama-2 dataset, and display the findings in \autoref{fig:entropy_ablation}. Our observations indicate a positive relationship between model scaling and both entropy model size and context window length; however, gains taper off when model capacity exceeds 50m parameters.

\paragraph{Types of Patching}

We conduct ablation studies on the four distinct patching strategies described in Section~\ref{section:patching}, namely: 1) Strided Patching employing strides of 4 and 6, 2) Patching according to whitespace, 3) Patching guided by BPE Tokenizer using the Llama 3 tokenizer, and 4) Entropy-driven patching utilizing a compact byte LLM.

Although dynamic patching shortens the actual sequence length, we standardize the sequence length across all patching strategies to ensure consistent context sizes. This approach guarantees that each model, on average, is exposed to the same quantity of bytes per sequence during both training and inference phases, thereby eliminating any biases related to larger context modeling. The outcomes of these ablation studies are presented in Figure~\ref{fig:scaling-compute-opt}. Notably, all alternative patching approaches surpass the performance of static patching, with space patching yielding results nearly equivalent to those of dynamic entropy-based patching.

In \autoref{tab:evals_ablation}, we report benchmark results for {\textsc BLT} models trained on the Llama 2 dataset, providing a comparative analysis of tokenizer-based models, space patching, and entropy-based patching, each optimized for the ideal number of training steps~\citep{dubey2024llama}. While space patching offers a more straightforward approach by avoiding real-time entropy model computations during training, our results confirm that the advantages seen with entropy-based patching in scaling experiments~(Section~\ref{section:scaling}) also extend to a variety of downstream benchmark tasks.\footnote{Results for space patching derive from earlier runs lacking cross-attention; however, the observed patterns persist even when cross-attention is incorporated.}

\paragraph{Cross-Attention} 
In~\autoref{tab:crossattn_ablations_1b}, we present an ablation study on the placement of cross-attention modules within the encoder and decoder components of {\textsc BLT}. For encoder-side cross-attention, our evaluation considers three approaches to initializing the queries: 1) employing an identical learned embedding across all global states, 2) utilizing a hash-based embedding derived from the patch bytes, and 3) aggregating the encoder's hidden states for the patch bytes at the corresponding layer via pooling.

Our results indicate that integrating cross-attention within the \textit{decoder} yields the strongest performance. Implementing cross-attention in the encoder brings modest benefits, but these are evident only when queries are initialized via pooling. Furthermore, our analysis reveals that cross-attention proves especially advantageous on Common-Crawl data, with the effect becoming more pronounced for larger patch sizes.

\paragraph{n-gram Hash Embeddings} 

We experiment with varying the n-gram hash embedding vocabulary sizes at 0, 100K, 200K, and 400K, and report the corresponding outcomes in Table~\ref{tab:ngram_ablations_1b}. Across all domains, we observe that incorporating hash embeddings yields benefits, with the most pronounced effects on Wikipedia and Github (showing a 0.04 bpb improvement versus a 0.01 bpb improvement after 15k steps at the 8B scale). Adjusting the hash vocabulary from 500K to 300K at the 8B model size led to a minimal performance change of approximately 0.001 bpb after 15k steps. This suggests that hash embeddings are crucial for enabling {\textsc BLT} to reach the performance level of models utilizing tokenizers, though enhancements diminish when surpassing 300K hash entries. Furthermore, our results suggest that these improvements are largely additive to those achieved via cross-attention, as they impact different datasets in complementary ways.

\paragraph{Local Model Hyperparameters} Table \ref{tab:local_layers} presents an ablation study across different configurations regarding the number of layers utilized in both the local encoder and decoder. Notably, when combined with hash n-gram embeddings, {\textsc BLT} demonstrates strong performance even when the encoder is minimal, employing only a single layer, provided that the decoder has a greater number of layers.

\section{Related Work}
\label{section:related_work}

\textbf{Character-Level RNNs:} The task of Character Language Modeling has maintained its significance in neural modeling research since its inception~\citep{sutskever2011generating,mikolov2012subword,graves2013generating}, largely due to its inherent capability to naturally handle out-of-vocabulary tokens, removing reliance on traditional back-off techniques. For instance, \cite{kim2016character} propose a framework that employs exclusively character-based input, utilizing convolutional and highway layers feeding into LSTM-based RNNs. This architecture achieves parity with then state-of-the-art RNN-driven language models on English, while yielding superior outcomes on morphologically complex languages, showcasing a key benefit of character-level LLMs. In a related direction, \cite{kenter2018byte} explore byte-level LSTM networks for machine comprehension tasks, showing improved performance over word-based models for highly inflected languages like Turkish and Russian. Similarly, \cite{zhang2015character} demonstrate that convolutional models built upon character representations can surpass word-level approaches in certain classification problems. Further advancements, such as those by \citet{chung2019hierarchical}, use hierarchical LSTMs with boundary-detecting modules at each tier to uncover latent structural hierarchies in text, advancing the effectiveness of character-level modeling. In the case of ByteNet~\cite{kalchbrenner2016neural}, CNN-based architectures applied to characters are proposed as an alternative to attention mechanisms in the context of machine translation.

\textbf{Character-Level Transformers:} The introduction of transformers utilizing attention mechanisms~\citep{vaswani2017attention} and subword tokenization techniques~\citep{sennrich-etal-2016-neural} has led to notable advances in language modeling and benchmark evaluations for neural architectures. Yet, segmenting text into word or sub-word representations imposes inherent inductive biases regarding the abstraction layers the models are exposed to. In an effort to leverage both the strengths of transformer architectures and the encouraging outcomes from character-level language modeling, \citet{al2019character} employed substantially deep transformer networks. By incorporating auxiliary loss functions, they managed to train transformer-based models that surpassed earlier LSTM-based character language models in performance. Nonetheless, these models continued to lag noticeably behind large language models operating at the word level. Similarly, GPT-2~\citep{radfordgpt2} documented that when trained on extensive corpora such as the 1 Billion Word Benchmark, language models based on byte-level representations did not match the efficacy of word-level approaches.

Although \cite{Choe2019BridgingTG} showed that transformer-based LLMs operating at the byte level can outperform comparable subword-level models in terms of accuracy, these byte-level models require significantly more computational resources and time to train. In a related vein, \cite{el2020characterbert} introduced a BERT variant called CharFormer, which constructs word representations via convolutional operations on character embeddings. They report performance gains in the medical domain, but again, at the cost of substantially higher computational expenditure. The work by \cite{clark2022canine} presented CANINE, an encoder-only model with 150M parameters that processes character sequences directly. CANINE features a deep transformer stack, functionally analogous to our global model, along with a combination of a local transformer and strided convolutions to perform input downsampling. On downstream multilingual tasks, CANINE surpasses an equivalent token-level model (mBERT) in performance. The ByT5 model~\citep{xue2022byt5} pursued fully byte-level encoder-decoder models without relying on patching or segmentation operations. ByT5 achieved superior robustness to input noise and matched or exceeded the performance of tokenizer-based models with only a quarter of the training data, but the absence of patching led to highly resource-intensive attention operations applied over every input byte. Modeling text at the byte level rather than the subword level drastically increases input sequence length, creating scalability challenges for efficient training. Most recently, leveraging the Mamba Architecture~\citep{gu2023mamba}—which sustains a fixed-size memory state over long sequences—\cite{wang2024mambabyte} developed a byte-level Mamba model. This approach also forgoes patching and, at a parameter count of 350M and in settings with controlled computational effort, outperforms byte-level transformer architectures on the bits-per-byte metric across several datasets.

\textbf{Patching-based approaches:} Leveraging patching strategies can substantially reduce the elevated flop count typical of byte-level LLMs while aiming to maintain their performance; this potential has been validated by several studies on relatively small models and limited training datasets. For instance, \cite{nawrot-etal-2022-hierarchical} examine static patching with both downsampling and upsampling, and introduce the hourglass transformer, which achieves superior results compared to existing byte-level baselines at the 150M parameter scale. Building on this, \cite{nawrot-etal-2023-efficient} extend the methodology by employing dynamic patching mechanisms, such as a boundary predictor trained either end-to-end or with guidance from certain tokenizers, alongside an entropy-based model similar to {\textsc BLT}. Their experiments demonstrate that this technique surpasses contemporary vanilla transformers in language modeling tasks at a scale of 40M parameters and roughly 400M tokens. Additionally, \cite{lester2024training} study the impact of training on data compressed via arithmetic coding, which enables them to attain greater compression rates than BPE. Using the equal-info windows approach, they achieve improved performance over byte-level models on language modeling tasks, though their approach still falls short of subword model baselines.

This study is heavily influenced by and most directly compares to MegaByte~\citep{yu2023megabyte}, a causal decoder-only large language model (LLM) that implements a fixed, static approach for segmenting and merging representations, thereby converting bytes into patches. On the decoder end, MegaByte utilizes a local model to revert from patches back into bytes. Their results show that MegaByte is capable of achieving performance on par with models using tokenizers, specifically at the 1B parameter level, when trained on a corpus of approximately 400B bytes. In all our experiments, we include MegaByte as a baseline, observing that its static patching technique consistently underperforms relative to current state-of-the-art tokenizer-based models, particularly when training is controlled for floating-point operations (FLOPs). We then illustrate how {\textsc BLT} effectively narrows this performance gap. Similarly, \cite{slagle2024spacebyte} recognize the limitations of MegaByte and propose an extension involving patching at whitespaces and analogous byte classes, in addition to integrating a local encoder module. Their research notes better performance compared to tokenizer-based transformer architectures within compute-constrained settings for certain datasets such as Github and arXiv, especially at the 1B parameter level. Our findings with this extended model support their conclusions, and we highlight that additional modifications in model architecture are essential to push byte-level approaches further and reach parity with leading token-based systems like Llama 3.

\section{Limitations and Future Work}

In this study, we adopt the number of training steps suggested as optimal for Llama~3~\citep{dubey2024llama} when making architectural decisions. Nevertheless, these scaling laws are based on transformers with BPE-level tokenization and might not yield ideal ratios of data to parameter sizes for {\textsc BLT}. Establishing scaling laws specifically tailored to {\textsc BLT}—which could potentially reveal more advantageous scaling patterns for this architecture—is reserved for future investigation. Furthermore, our experiments primarily focus on model sizes up to 1B parameters; thus, it remains possible that scaling up to 8B parameters and beyond could necessitate different architectural strategies, potentially resulting in superior performance at larger model sizes.

Current transformer libraries and codebases prioritize efficiency when working with tokenizer-based transformer models. Although we provide theoretical flop-equivalent experiments and employ select efficient solutions like FlexAttention for layers that diverge from the standard transformer architecture, our current implementations may still lag behind tokenizer-based models in actual wall-clock performance, indicating room for additional optimization.

Whereas {\textsc BLT} employs an independently trained entropy model for patching, an alternative avenue worth investigating is to jointly train the patching model in an end-to-end manner. In Section \ref{sec:distillation}, we provide preliminary experimental evidence suggesting that it is feasible to adapt tokenizer-based models—such as Llama 3, which has been trained on over 10T tokens—into byte-level models by initializing and freezing the global transformer with the pretrained weights. Pursuing this line of research could reveal approaches that preserve the advantages of bytefying while potentially achieving superior performance compared to tokenizer-based models, eliminating the need for training entirely from scratch.

\section{Conclusion}
\label{section:conclusion}

In this work, we introduce the Byte Latent Transformer (\textbf{BLT}), an innovative model architecture that challenges the prevailing reliance on fixed-vocabulary tokenization in large-scale language modeling. {\textsc BLT} incorporates a dynamic and trainable mechanism for assembling bytes into patches, dynamically tailoring computational expenditure to the intrinsic complexity of the input data. This adaptive strategy yields marked advances in both computational efficiency and model robustness. Our comprehensive scaling experiments show that {\textsc BLT} achieves performance on par with established tokenization-dependent architectures, such as Llama 3, when scaled up to 8B parameters and 4T bytes of data, and achieves as much as a 50\% reduction in inference floating point operations with only marginal degradation in evaluation scores. Additionally, {\textsc BLT} introduces a novel approach to scaling: it enables concurrent growth of both model parameters and patch sizes, even within a fixed inference compute budget—a feature that is especially beneficial under typical compute limitations encountered in real-world applications. By processing raw byte-level data directly, {\textsc BLT} enhances a model’s capacity for managing infrequent or atypical input sequences, bolstering both its resilience to noisy data and its sensitivity to sub-word phenomena. Altogether, these findings establish {\textsc BLT} as a viable and compelling substitute for conventional tokenization-centric methods, offering a flexible and efficient framework for the next generation of language models.

\end{document}